\documentclass[a4paper,12pt]{article}

\usepackage[utf8]{inputenc}
\usepackage[T1]{fontenc}
\usepackage[brazil]{babel}
\usepackage{fancyhdr}
\usepackage{graphicx}
\usepackage{float}
\pagestyle{fancy}
\setlength{\headheight}{14.49998pt}
\fancyhf{}

\title{Game Design Document}
\author{
    Henrique Finatti Silveira Belo Trebbi \\
    Tiago Fagundes dos Santos \\
    Mateus Marana Assuena \\
    Giovanni Chahin Morassi \\
}
\date{Março de 2026}

\renewcommand{\headrulewidth}{0.5pt}
\fancyhead[L]{FeiKemon}
\fancyhead[R]{CC7140 - Jogos Digitais}

\renewcommand{\footrulewidth}{0.5pt}
\fancyfoot[L]{FeiKemon}
\fancyfoot[C]{\thepage}
\fancyfoot[R]{FEI - SBC}

\begin{document}

\maketitle
\newpage
\section{FeiKemon}
\subsection{Sistemas de jogo}
\subsubsection{Movimentação no mapa}
O jogador poderá explorar o mapa utilizando os botões de movimentação “A”, “W”, “S” e “D”.

\subsubsection{Interação com os NPCs}
O jogador poderá interagir com os NPCs do mapa chegando perto o bastante e clicando com a barra de espaço.
Um som será emitido quando essa interação ocorrer.

\subsubsection{Combate}
Quando uma batalha começa, seja contra um NPC ou com um FeiKemon aleatório,
o jogador pode utilizar os botões de movimento para selecionar o tipo de interação.
“Atacar”, “Fugir” ou “Capturar” ou, caso tenha selecionado “Atacar”,
o jogador pode selecionar os golpes com os botões de movimento também.

\subsection{Faixa etária recomendada}
A faixa etária recomendada é para pessoas de pelo menos 17 anos.

\subsection{Classificação indicativa ESRB pretendida}
T (Teen). Pode conter violência, temas sugestivos, humor grosseiro, mínimo sangue e uso ocasional de linguagem forte.

\subsection{Resumo da história do jogo}
O jogo é um jogo top-down, que se passa no dia da recepção dos calouros da faculdade FEI, e tem como personagem principal
um ansioso garoto que acaba de se matricular na sua faculdade.

Durante a recepção o aluno terá uma palestra da reitoria onde irá conhecer algumas curiosidades sobre a faculdade e os FeiKemons,
após isso será convidado a conhecer a faculdade e os prédios dos cursos disponíveis. Durante sua exploração o jogador irá perceber que
há algo estranho com a faculdade e que só ele e o professor Fagner percebem isso, o aluno irá mais a fundo na história para descobrir que a faculdade
Mauá está impactando negativamente a faculdade e ele deve salva-la.

Para salva-la o aluno irá montar um time de FeiKemons e derrotar os mestres da FEI para remover os efeitos da Mauá, mas para isso ele precisa primeiro derrotar
a chefe de departamento de computação para conseguir acesso a sala dos mestres.
Os mestres da FEI são a reitoria, o padre e professor Samir.

\subsection{Modos de jogo distintos}
Modo História

\subsection{Diferenciais competitivos}
O jogo se diferencia por ter foco na integração do verdadeiro novo aluno com uma versão mais amigável e animada das suas novas experiências na faculdade. O aluno consegue relacionar a experiência que ele tem na vida real com a experiência que ele tem jogando o jogo.

\subsection{Produtos concorrentes}
Jogos da franquia Pokemon.
\end{document}